\documentclass[aspectratio=169]{beamer} % 16:9比例
% 设置字体支持
\usepackage[UTF8]{ctex}
\usepackage[T1]{fontenc}
\usefonttheme[onlymath]{serif} % 数学公式使用衬线字体
\setsansfont{Arial} % 设置无衬线字体为Arial

% 引入常用宏包
\usepackage{multicol} % 用于显示多栏内容
\usepackage{booktabs} % 用于美化表格
\usepackage{calligra} % 用于感谢页的手写字体
\usepackage{graphicx} % 用于插入图片
\usepackage{fontawesome5}

\usepackage{zju} % 自定义的beamer主题样式文件
\usepackage{ragged2e} % 提供了更好的两端对齐算法

\title{浙江大学Beamer模板}
\author{徐翌森}
\institute{系统5班}
\date{\today}

\begin{document}

\begin{frame}[plain,noframenumbering]
    \titlepage
    \vspace{-0.5cm}
    \begin{figure}[h]
        \centering
        \includegraphics[width=0.15\textwidth]{zju/zju_logo.pdf}
    \end{figure}
\end{frame}

\section{第一章节}

\begin{frame}{第一章内容}
    \textcolor{zju_blue}{蓝色文字}
\end{frame}

\section{第二章节}

\begin{frame}{第一章内容}
    \textcolor{zju_red}{红色文字}
\end{frame}

\appendix

\section*{致谢}

\begin{frame}[plain,noframenumbering]
    \begin{center}
        \vspace{1cm}
        \hrulefill 
        \raisebox{-2ex}{\quad {\Huge\calligra Thanks!} \quad} 
        \hrulefill

        \vspace{1cm}

        {\Large 请老师批评指正}
    \end{center}
\end{frame} 

\end{document}